\documentclass[11pt,letterpaper]{article}
\usepackage[T1]{fontenc}
\usepackage[utf8]{inputenc}
\usepackage{lmodern,microtype}
\usepackage[margin=0.78in,headheight=15pt]{geometry}
\usepackage{amsmath,amssymb,booktabs,longtable,tabularx,array,enumitem,xcolor,hyperref,fancyhdr,lastpage,needspace}
\definecolor{EngineBlue}{HTML}{183B56}
\definecolor{EngineTeal}{HTML}{117A83}
\definecolor{EngineGray}{HTML}{596671}
\definecolor{EngineLight}{HTML}{EEF4F7}
\hypersetup{colorlinks=true,urlcolor=EngineTeal,linkcolor=EngineBlue}
\setlength{\parindent}{0pt}
\setlength{\parskip}{5pt}
\emergencystretch=3em
\setlist{nosep,leftmargin=*}
\pagestyle{fancy}\fancyhf{}
\fancyhead[L]{\small\sffamily Game Engine Architecture | Engineering Suite}
\fancyhead[R]{\small\sffamily \DocType}
\fancyfoot[L]{\scriptsize Source-grounded engineering curriculum | v1.0}
\fancyfoot[R]{\scriptsize \thepage/\pageref{LastPage}}
\newcommand{\DocType}{Reference}
\newcommand{\SuiteHeading}[3]{%
\renewcommand{\DocType}{#2}%
{\small\bfseries\sffamily\color{EngineTeal}#2\hfill #3}\par\vspace{6pt}
{\LARGE\bfseries\color{EngineBlue}#1}\par\vspace{4pt}
{\color{EngineTeal}\rule{\linewidth}{1pt}}\par\vspace{7pt}}
\newcommand{\LabEvidence}[1]{\par\textbf{Evidence required:} #1\par}
\newcommand{\WriteLines}[1]{\par\begingroup\setlength{\parskip}{9pt}\count255=0\loop\ifnum\count255<#1\noindent\rule{\linewidth}{0.35pt}\par\advance\count255 by1\repeat\endgroup\smallskip}
\newcommand{\WritingBox}[1]{\noindent\fbox{\parbox[t][#1][t]{0.97\linewidth}{\mbox{}}}\par\smallskip}
\newcommand{\SourceNote}{\footnotesize\textit{Primary reading: Jason Gregory, \emph{Game Engine Architecture}, 3rd ed. (CRC Press, 2018). The narrative, exercises, laboratory designs and scoring instruments in this suite are independently authored teaching materials, not a reproduction of the textbook.}\normalsize\par}
\newcommand{\Criterion}[2]{\textbf{#1.} #2\par}
\newcommand{\FormField}[2]{\Needspace{3\baselineskip}\textbf{#1}\par\textcolor{EngineGray}{\footnotesize #2}\par\WriteLines{2}}
\newcommand{\BlankField}[1]{\Needspace{3\baselineskip}\textbf{#1}\par\WriteLines{2}}
\newcolumntype{Y}{>{\raggedright\arraybackslash}X}

\begin{document}
\SuiteHeading{Chapter 17 — You Mean There's More?}{Study guide}{Part Conclusion}
\SourceNote
\section{Reading crosswalk and prerequisites}
Book chapter 17; Part Conclusion; sections 17.1--17.2. Read the source chapter first, then perform the independent engineering activities below. Distinguish statements from the book from proposed implementations and experiments.

Required prior evidence: selected exercises from chapters 15 and 16, or a documented equivalent baseline.

\section{Observable learning outcomes}
\begin{itemize}
\item Explain and defend system boundaries in terms of contracts, ownership and measurable behavior.
\item Explain and defend unsupported engine topics in terms of contracts, ownership and measurable behavior.
\item Explain and defend roadmap prioritization in terms of contracts, ownership and measurable behavior.
\item Explain and defend integration risk in terms of contracts, ownership and measurable behavior.
\end{itemize}
\section{Engineering interpretation}
A learning engine is deliberately incomplete. Architecture maturity is measured by verified requirements, explicit omissions and well-supported trade-offs.

\section{Section-by-section reading tasks}
Read each named section in the source and complete the corresponding original engineering prompt. Keep a source page reference beside every claim.
\small\begin{longtable}{|p{0.10\linewidth}|p{0.25\linewidth}|p{0.55\linewidth}|}
\hline \textbf{Sec.} & \textbf{Book topic} & \textbf{Independent engineering task} \\ \hline\endhead
17.1 & Some Engine Systems We Did Not Cover & Register gaps such as networking, UI, online services, authoring tools and localization. \\ \hline
17.2 & Gameplay Systems & Define future gameplay extension points and stable ownership boundaries. \\ \hline
\end{longtable}\normalsize
\section{Guided reading and reasoning}
\begin{enumerate}
\item Inventory subsystems absent from the learning engine (e.g., networking, localization, advanced AI) rather than claiming a complete production engine.
\item Prioritize work by player requirements, team capacity and platform constraints.
\item Defend one build-vs-buy decision using non-functional quality attributes.
\end{enumerate}
\section{Implementation exercise}
Write a staged extension roadmap identifying interface dependencies and test gates.

\subsection{Minimum engineering specification}
\begin{itemize}
\item State all preconditions, postconditions and owned resources in a short interface contract.
\item Make a deterministic test case and a negative/error-path test case.
\item Record CPU, platform, compiler or Python version, operating system, seed, and relevant runtime settings.
\item Keep the exercise separate from proprietary game code or assets; produce original diagrams and sample data.
\end{itemize}
\section{Validation and adversarial cases}
Run a final architecture review and record material risks and deferred capabilities.

\begin{itemize}
\item What invariant could fail while a demo still appears correct?
\item What input or platform assumption invalidates your design?
\item How would you diagnose a regression from logs, measurements and a minimal reproduction?
\end{itemize}
\section{Study questions}
\begin{enumerate}
\item What is the principal abstraction boundary introduced in this chapter?
\item Which performance/correctness trade-off does that boundary enable?
\item How does this subsystem coordinate with resource lifetime, simulation time and threading?
\item What would be a counterexample to your assumed invariant?
\end{enumerate}
\section{Exit criteria and artifacts}
\begin{itemize}
\item One-page technical decision note with alternatives and constraints.
\item Runnable source or a precise pseudocode specification; include test commands.
\item Test transcript, benchmark/data capture where applicable, and a written conclusion.
\item Self-assessment: identify one unresolved hazard and how you would validate it.
\end{itemize}

\end{document}
